\section{1 Paskaita}

\subsection*{Pagrindinės atomų savybės}

Atomas neturi apibrėžtos ribos, nes jo sandarą lemianti elektrostatinė jėga yra toliasiekė

\begin{itemize}
    \item Mažiausias matomas dalykas $10^{-6}m$\
    \item Atomas $10^{-6}m$
    \item Branduolys $10^{-15}m$
\end{itemize}

Atomo matmenys apribūdina:
\begin{itemize}
    \item Boro orbitos spindulys $a$
    \item Vidutinis išorinio elektrono atstumas nuo branduolio
\end{itemize}

Vandenilio atome:
\begin{itemize}
    \item Pirmosios Boro orbitos spindulys
        \begin{equation*}
            a_0 = \frac{\hbar}{m_e e^2} \simeq 0.5 \cdot 10^{-10}m
        \end{equation*}
    \item $n$-tosios Boro orbitos spindulys
        \begin{equation*}
            a_{0n} = n^2 a_0 
        \end{equation*}
\end{itemize}

Vidutinis elektrono atstumas nuo branduolio:

\begin{equation*}
    \bar{r} = \int \Psi*(\bm{r}) \bm{r} \Psi{\bm{r}} d \bm{r} \qquad \bar{r}_{nl} ~ \frac{1}{Z} 
\end{equation*}
čia $nl - 1s,2s,2p,3s,3p,3d,4s$. $n$ - pagrindinis kvantinis skaičius (1,2,3,...), $l$ - orbitinis kvantinis skaičius (s,p,d,...)

Atomo masė: $1u \simeq 1.66057 \cdot 10^{-27} kg$. Atominis masės vienetas - tai anglies izotopo ${}^{12}C$ atomo massės $\frac{1}{12}$ dalis.

Atomo branduolio dydis: $R = (1.3 \pm 0.1)A^{\frac{1}{3}}10^{-15}m$, čia $A$ - masės skaičius (sveikas skaičius, artimiausias santykinei atomo masei.) Pvz
\begin{itemize}
    \item $R({}^{12}C)=2.7 \cdot 10^{-15} m$
    \item $R({}^{208}Pb)=7.10 \cdot 10^{-15} m$
\end{itemize}

Teigiamas jonas - atomas, netekęs elektronų
Neigiamas jonas - atomas, prisijungęs elektronus
Egzotinis atomas - prisijungęs pozitroną, miuoną arba kitokią elementariąją dalelę

Atomas, netekęs $n$ elektronų - $n$-kartinis jonas $A^{n+}$. Pvz $Ar^{4+} \equiv ArV$ - keturkartinis jonas.

Miuoniniai atomai - gaunami dirbtiniu būdu: protonais bomborduojant medžiagą; kai kurie sukurti miuonai yra pagribiami atomų. Miuonas - sunkusis elektronas, toks pat krūvis, 207 kartus didesnė masė. Miuonas nestabilus, suyra į kitas daleles per $10^{-6}s$. Dėl didelės masės miuonas sukasi labai arti branduolio. Jo spinduliutė labai priklauso nuo branduolio savybių - krūvio pasiskirstymo jame, formos, magnetinio momento. Miuoninių atomų spektrai teikia svarbią informaciją apie atomų branduolius.

Jonų susidarymas gamtoje ir laboratorijose. Žvaigždžių spektruose stebimos nuo kelių iki keliolikos jonizuotų atomų spektro linijos. Laboratorijose galima gauti bet kurio elemento, bet kurį joną.

EBIT - electron beam ion trap.

Atomas ir Saulės sistema:
\begin{itemize}
    \item Panašumai:
    \begin{itemize}
        \item Mažas masyvus centras
        \item Aplink skrieja mažesnės masės kūnai
    \end{itemize}
    \item Skirtumai:
    \begin{itemize}
        \item Atomas - sfera, Saulės sistema - diskas,
        \item stiprios elektrinės jėgos ir silna gravitacijos jėga
        \item elektronai identiški, planetos skirtingos
        \item Saulės sistemoje veikia klasikinė mechanika, tikslio padėtys, o atome - kvantinė mechanika, tikimybinis pasiskirstymas.
    \end{itemize}
\end{itemize}

Pagrindiniai dydžiai:
\begin{itemize}
    \item $a_0 = 0.52918 \cdot 10^{-10} m$ - pirmosios Boro orbitos spindulys
    \item $m_e=9.1095 \cdot 10^{-10} kg$
    \item $\hbar = \frac{h}{2 \pi} = 1.0546 \cdot 10^{-34} Js$ - veikimo vienetas
    \item $e=1.6022 \cdot 10^{-19} C$
\end{itemize}

Išvestiniai dydžiai:
\begin{itemize}
    \item Energijos vienetas 
    $\frac{m_e e^4}{\hbar^2} = 4.3438 \cdot 10^{-18} J$
    
    $1eV = 1.6022 \cdot 10^{-19}J$

    $1 a.v. = 27.213 eV$

    \item Laiko vienetas:
    $\frac{\hbar^3}{m_e e^4} = 2.4192 \cdot 10^{-16} s$

    \item Atominis greičio vienetas
    $\frac{e^2}{\hbar} = 2.1877 \cdot 10^6 m/s$

    $c \simeq 137 a.v.$
\end{itemize}


Lygtys:
\begin{itemize}
    \item Klasikinėje mechanikoje - II Niutono dėsnis
    \item Kvantinė mechanioje - Schrodingerio lygtis
\end{itemize}

Stacionarioji (nekintanti laike) Schrodingerio lygtis:

\begin{equation*}
    \hat{H} \Psi = E \Psi
\end{equation*}
čia $\hat{H}$ yra energijos operatorius Hamiltonianas, $E$ energija, $\Psi$ sistemos tikrin4 banginė funckija. Vieno laisvo atomo atveju paprastai priimamos prielaidos:
\begin{itemize}
    \item Atomo branduolys nejuda (su juo susijama koordinačių sistemos pradžia)
    \item Branduolys yra taškinis
\end{itemize}

%-----------------

Lygtis radialijai funkcijai gaunama, įstačius į vienelektronę Schrodingerio lygtį:

\begin{equation*}
    \hat{h} \psi_{nlm} = \epsilon_{nl} \psi_{nlm}
\end{equation*}

Hamiltoniano išraišką
\begin{equation*}
    \hat{h} ž \frac{p^2}{2m_e}
\end{equation*}


%------------------

Pertvarkius gaunasi tokia lygtis radialiajai funcijai (atominiais vientetais):
\begin{equation*}
    \left( \frac{d^2}{dr^2} - 2 V_{nl}(r) - \frac{l (l+1)}{r^2} \right) P_{nl}(r) = -2 \epsilon_{nl} P_{nl}(r)
\end{equation*}

Sprendinys turi tenkinti kraštines sąlygas: $P_{nl} (0) = 0$, $P_{nl}(\infty) = 0$.

Sprendiniai egzistuoja tik tam tikroms tikrinėms $\epsilon_{nl}$ vertėms, sudarančioms diskretinį sprektrą. Skaičius $n$ numeruoja lygties sprendinius, pradedant nuo $n=l+1$.

%----------------


